\section{Effective front roots for all minimum cycle lengths}

\textbf{Partial theorem.} Let $n\ge10{,}000$ and $3\le k\le n$. Set $H=\lfloor7\sqrt n/20\rfloor$ and $P_n=E_{n,k}/t$. For every $1\le h\le H$, the root of $P_n$ of absolute rank $h$ is simple and negative. At each such root $x$, the distinguished-cycle term $L=\binom n{k-1}A_{k-1}E_{n+1-k,k}$ in the insertion recurrence satisfies
\[
|L(x)|<|x(1-x)P_n'(x)|.
\]
In order of increasing absolute value, the new rank-$h$ root is strictly closer to zero than the old rank-$h$ root. All higher-rank new roots are farther from zero than that old root. All $H$ old roots precede the phase sample $t=-\exp(-\pi\cot((H+1)\pi/n))$. The reciprocal conclusions hold at the other end. There is no restriction on the number of cycles. The proof compares the run-color probabilities for consecutive rows and bounds the remaining terms on three circles.

Use the \href{https://github.com/ProgramComputer/open-math-research/blob/main/topics/long-cycle-excedances/proof.md}{run encoding and insertion recurrence}. For each $p\ge1$ write $\alpha_{N,p}=[t^p]E_{N,k}/p^N$, which is the probability that an independent uniform word of length $N$ on $p$ run colors and a uniform permutation of the colors satisfy all upward boundaries and minimum lengths. Write $f(N,p,k)$ for the acceptance probability for the minimum lengths alone. It bounds $\alpha_{N,p}$ above. Put $f(N,p,k)=0$ for $0<N<k$ when $p\ge1$. The good event that every color appears in both label halves guarantees the upward boundaries and all nonempty colors, so
\[
0\le f(N,p,k)-\alpha_{N,p}\le 2p e^{-\lfloor N/2\rfloor/p}.
\tag{F1}
\]
This argument uses the probability form of the coefficients from the \href{https://github.com/ProgramComputer/open-math-research/blob/main/topics/long-cycle-excedances/uniform-interlacing.md}{uniform proof} and the adaptive root scaling from the \href{https://github.com/ProgramComputer/open-math-research/blob/main/topics/long-cycle-excedances/two-cycle-front.md}{two-cycle front estimate}.

\subsection{Effective run-color acceptance envelopes}

Put $J=\lfloor\sqrt n/2\rfloor$. For $p\le J$ and $a=kp/n$, the old row satisfies
\[
\frac{99}{100(1+(8/3)a)}\le\alpha_{n,p}\le
\min\left(1,\frac{151}{100a}\right).
\tag{F2}
\]
The upper bound also applies to $f(N,p,k)$ with $a$ replaced by $kp/N$ for every $n/2\le N\le n$ and $1\le p\le J$.

For the lower bound when $k\le n/2$, let $b=\lceil (4/3)a\rceil$, so $b\le p$. Write $d(p,b)$ for the probability that a uniform permutation of $p$ colors has no cycle shorter than $b$. Distinguishing the cycle containing a fixed color gives $1/(2b-1)\le d(p,b)\le1/b$ for $p\ge b$. Condition the run permutation to have no cycle shorter than $b$. A cycle of $u\ge b$ colors receives a binomial number of labels of mean $nu/p\ge4k/3$, so its probability of receiving fewer than $k$ labels is at most $e^{-nu/(32p)}$. For $u<p$, the unconditional expected count of marked $u$-cycles whose remaining cycles also have length at least $b$ is $d(p-u,b)/u\le1/(bu)$. A cycle of length $p$ receives every label and cannot fail. After division by $d(p,b)$, the union bound over the marked cycles is therefore less than
\[
2\sum_{u\ge1} \frac{e^{-nu/(32p)}}{u}
\le \frac{2e^{-25/4}}{1-e^{-25/4}}<\frac1{250}.
\]
The error in (F1) relative to $d(p,b)$ is bounded by $4p^2 e^{-(n-1)/(2p)}\le n e^{-(n-1)/\sqrt n}<1/1000$. Thus $\alpha_{n,p}\ge(99/100)d(p,b)\ge99/[100(1+(8/3)a)]$. When $k>n/2$, acceptance for the size condition requires a single run cycle and has probability $1/p$. Applying (F1) gives the same lower bound.

For the upper bound take $\mu=N/p\ge\sqrt n$, $a'=k/\mu$ and $b=\max(1,\lceil (2/3)a'\rceil)$. The bound is immediate if $k>N$ or $a'\le3/2$. Otherwise, acceptance with a run cycle of $u<b$ colors requires that this cycle receive at least $k$ labels despite mean $Nu/p<(2/3)k$. The exponential binomial tail is at most
\[
\exp[-k(\log(3/2)-1/3)]\le e^{-k/15}.
\]
Using the unconditional expected number $1/u$ of $u$-cycles gives
\[
f(N,p,k)\le d(p,b)+(1+\log p)e^{-k/15}
\le\frac{3}{2a'}+(1+\log p)e^{-a'\mu/15}.
\]
The last error times $a'$ is at most $(3/2)(1+\log(\sqrt n/2))e^{-\sqrt n/10}<1/100$. This uses decrease in $a'$ for $a'\ge3/2$, decrease in $\sqrt n$ for $\sqrt n\ge100$, and the bounds $p\le\sqrt n/2$ and $\mu\ge\sqrt n$. This proves (F2).

\subsection{Lag coefficients and consecutive row stability}

Write $m=n+1-k$. If $m<k$ the lag vanishes and we set $\lambda_p=0$. Otherwise $n/2\le m\le n$, and for $p\ge1$ we put
\[
\lambda_p=\sum_{j=1}^{p-1} b(j/p)f(m,p-j,k),
\qquad
b(y)=\binom n{k-1}y^{k-1}(1-y)^m.
\]
Then
\[
[t^p](tL)\le p^n\lambda_p,\qquad
0\le f(n+1,p,k)-f(n,p,k)\le \lambda_p/p .
\tag{F3}
\]
For the first bound, the coefficient of $t^j$ in $tA_{k-1}$ is at most $j^{k-1}$ by assigning the $k-1$ labels of a permutation to its increasing runs. Convolve with $[t^{p-j}]E_{m,k}=(p-j)^m\alpha_{m,p-j}$. Division by $p^n$ uses $k-1+m=n$, giving the one binomial point mass in (F3). For the second bound, couple the word of length $n+1$ to the old word by appending an independent color. If the minimum-size condition changes from false to true, the new label hits a run cycle that had precisely $k-1$ labels. A cycle of $j$ colors contributes expected count $1/j$ and is hit with probability $j/p$. Conditional on the count $k-1$, the remaining labels use the other $p-j$ colors independently and uniformly, and the remaining run permutation is uniform. The all-color cycle cannot have had size $k-1$. The same argument makes both differences zero when $m<k$, and both vanish for $p=1$. For use outside the range of (F2), $\lambda_p\le p$ follows directly.

For $p\le J$ the useful uniform estimate is
\[
\lambda_p\le \frac8{25}\alpha_{n,p}.
\tag{F4}
\]
If $k\le3n/(4p)$, set $\mu=n/p$. The sample $b(j/p)$ is bounded by $\Pr\{\mathrm{Bin}(n,j/p)\le k-1\}\le e^{-j\mu/32}$ for every $j\ge1$. Hence $\lambda_p<1/500$, while $\alpha_{n,p}\ge33/100$ by (F2).

Otherwise $k-1\ge150$. Put $y_0=(k-1)/n$. A sample with $j/p\le y_0+1/10$ has residual parameter
\[
\frac{k(p-j)}m\ge\frac45\frac{kp}n
\]
because $m/n=1-y_0\ge1/2$. The upper bound from (F2) for the residual and the lower bound for $\alpha_{n,p}$ give
\[
\frac{f(m,p-j,k)}{\alpha_{n,p}}
\le \frac{100}{99}(1+(8/3)a)\min\left(1,\frac{151}{80a}\right)
<\frac{61}{10}.
\]
For $j/p>y_0+1/10$, the ratio of $b(j/p)$ to its maximum at $y_0$ is $\exp(-nD(y_0\|j/p))$. Here $D$ is Bernoulli relative entropy and $D(v\|w)\ge2(v-w)^2$, by the second derivative in $v$. Hence $b(j/p)\le e^{-n/50}$. Summing these samples after division by $\alpha_{n,p}$ costs at most $n e^{-n/50}<1/10{,}000$, using $p(1+(8/3)a)/(99/100)<n$.

Stirling's bounds give $b(y_0)\le e^{1/(12n)}/\sqrt{2\pi(k-1)(1-y_0)}<47/1000$. The lattice sum of the unimodal kernel is at most its integral times $p$, plus its maximum. This follows by omitting the largest sample and charging each other sample to the adjacent interval toward the largest one. The beta integral is $1/(n+1)$. Since $p/(n+1)<1/200$, the total relative cost is at most
\[
\frac{61}{10}\left(\frac{47}{1000}+\frac1{200}\right)+\frac1{10{,}000}
<\frac8{25}.
\]
This proves (F4).

Equations (F1), (F3) and (F4) also give
\[
\frac{|\alpha_{n+1,p}-\alpha_{n,p}|}{\alpha_{n,p}}\le \frac{33}{100p}
\quad(1\le p\le J).
\tag{F5}
\]
Indeed the two boundary-error terms in (F1), multiplied by $p/\alpha_{n,p}$, cost at most $2 n^{3/2}e^{-(n-1)/\sqrt n}<1/100$ by (F2). The same upper bound for the two row sizes follows by using $\lfloor n/2\rfloor$ in both boundary terms. For $p=1$ the ratio between the two coefficients $\alpha$ is exactly one.

\subsection{Isolation at the adaptive radius}

Fix $1\le h\le H$. Put $\rho=\alpha_{n,h}/\alpha_{n,h+1}$, $u=\rho(h/(h+1))^n$ and
\[
F(z)=\frac{E_{n,k}(uz)}{\alpha_{n,h}h^n(uz)^h},
\qquad
b_p^*=\left(\frac ph\left(\frac h{h+1}\right)^{p-h}\right)^n.
\]
The terms at $p=h,h+1$ are exactly $1,z$. On either side of these terms the ratios of consecutive $b_p^*$ away from the central pair are at most $q=e^{-n/(h+1)^2}$. This follows by evaluating the largest ratios at the two immediate neighbors, where the exact bases are $1-1/h^2$ and $1-1/(h+1)^2$. The left tail is empty for $h=1$.

From (F2), for $p,v\le J$,
\[
\frac{\alpha_{n,p}}{\alpha_{n,v}}
\le \frac{51}{10}\max(1,v/p).
\]
Thus both $\rho$ and $\rho^{-1}$ are below 6 if $h\ge6$, and below $51/5$ for $h\le5$. In the first case $h+1\le(9/25)\sqrt n$, so $q<1/2200$. In the second case $q\le e^{-10000/36}$, which is more than enough for $q\max(\rho,\rho^{-1})<6/2200$ again. At $p=h-\ell$ the coefficient ratio to $\alpha_{n,h}$ is at most $(51/10)(\ell+1)$. At $p=h+1+\ell\le J$ the ratio to $\alpha_{n,h+1}$ is at most $51/10$. When $p=h+1+\ell>J$, we have $\ell>(7/50)\sqrt n$ and $h+1<(18/7)\ell$. Hence
\[
1/\alpha_{n,h+1}\le\frac{100}{99}(1+(8/3)(h+1))<8(\ell+1).
\]
These bounds, and $\alpha_{n,p}\le1$, give right-tail ratio at most $8(\ell+1)$ for every $p=h+1+\ell$.

Let $s$ be any upper bound on $q\max(|\rho z|,|\rho z|^{-1})$. At left index $h-\ell$ the absolute term is at most $(51/10)(\ell+1)s^\ell$, and at right index $h+1+\ell$ it is at most $8|z|(\ell+1)s^\ell$. Thus one can use $s=1/300$ on $|z+1|=1/8$, $s=3/500$ on $|z|=1/2$ and $s=9/1000$ on $|z|=3$. On these circles respectively the error in $F-(1+z)$ is less than $1/10$, $1/2$ and 2. These follow by multiplying $(51/10)+8|z|$ by $\sum_{\ell\ge1}(\ell+1)s^\ell=(1-s)^{-2}-1$. Rouché's theorem for $z^{h-1}F(z)$ gives respectively one root in the disk about $-1$, $h-1$ roots inside the half-radius circle and $h$ inside the radius-three circle. The disk sits in the resulting annulus. Its unique root is simple and negative by conjugation, and has absolute rank $h$.

At this root write $x=uz$, $X=|x|$. We have $|z+1|<1/8$. Since $F(z)=0$, the identity $xE_{n,k}'(x)/(\alpha_{n,h}h^n x^h)=zF'(z)-F(z)$ gives the two central terms $-1,0$. The remaining terms are bounded by
\[
\frac{51}{10}\left(\frac{1+s}{(1-s)^3}-1\right)
  +9\frac{2s}{(1-s)^3}<\frac3{20}\qquad(s=1/300).
\]
The two weights are $(\ell+1)^2$ on the left and $\ell(\ell+1)$ on the right, including the earlier coefficient-ratio factors. Thus $|xE'_{n,k}(x)|>(17/20)\alpha_{n,h}h^nX^h$.

\subsection{Signs and the join}

In the scale $\alpha_{n,h}h^n X^h$, (F3) and (F4) bound the contribution of indices $p\le J$ to $|xL(x)|$ by at most $(8/25)$ times the absolute sum of the terms of $F(z)$, which is less than $(8/25)(1+9/8+1/10)$. For $p>J$ use $\lambda_p\le p$, $\ell=p-h-1\ge14$, $p<4\ell$, and the bound $1/\alpha_{n,h+1}<8(\ell+1)$. Their normalized sum is less than $36\sum_{\ell\ge14}\ell(\ell+1)(1/300)^\ell<1/1000$. Therefore $|xL(x)|<(18/25)\alpha_{n,h}h^nX^h<|xE'_{n,k}(x)|$. At an old root $xP_n'(x)=E'_{n,k}(x)$, so $|L(x)|<|x(1-x)P_n'(x)|$. The latter term has sign $(-1)^h$ because all $h-1$ smaller roots are negative and simple, and $P_n(0)=1$. Both signed insertions preserve this sign.

For consecutive order, let $u_{n,h}$ and $u_{n+1,h}$ be the two adaptive radii. At $p=h,h+1$, (F5) bounds $\alpha_{n+1,p}/\alpha_{n,p}$ between $83/100$ and $117/100$ for $p\ge2$. At $p=1$ it equals one. The old root is therefore in the new rank-$h$ annulus $(u_{n+1,h}/2,3u_{n+1,h})$ in modulus. For $h\ge2$, its modulus divided by $u_{n+1,h}$ is between $(7/8)(83/117)$ and $(9/8)(117/83)(3/2)$. At $h=1$, it is between $(7/8)(83/100)2$ and $(9/8)(117/100)2$. The new polynomial has exactly one root in that annulus and sign $(-1)^{h-1}$ before that root. Its sign $(-1)^h$ at the old root places the old root strictly after it in modulus, with every other larger new root outside the annulus.

Finally every old front root has $X<(9/8)(51/5)(H/(H+1))^n$. At the next integer phase sample $\theta=(H+1)\pi/n$, one has $\pi\cot\theta<n/(H+1)$ and
\[
n\log(1+1/H)-\frac n{H+1}
\ge\frac n{2(H+1)^2}
\ge\frac1{2(9/25)^2}
>\log((9/8)(51/5)).
\]
Hence each of the $H$ front roots comes before this sample. The sample satisfies $n\sin^2\theta>6/5$ by $H+1>7\sqrt n/20$, $\sin y\ge y-y^3/6$ and $n\ge10{,}000$. Reflection gives the corresponding order and insertion signs at large modulus. The estimate leaves the middle-root argument as the remaining analytic step for this cutoff.

\section{Partial fractions for any number of occupied Eulerian poles}

This estimate compares the middle roots of consecutive rows with arbitrary numbers of permitted cycles. Let $S=\sin\theta$, $nS^2\ge1$, $n\ge100$, $n\ge k\ge32$. Put $D_0=\sqrt{n/k}$ and $U=3+\log(3D_0)$. Assume
\[
kS\ge2000 U. \tag{PA}
\]
We prove that the integer-phase signs and both insertion signs at old roots hold under (PA), within the residual-index range below. It suffices below to consider $n\ge2k-1$, since the earlier transitions are single-cycle Eulerian. For the lag estimates below first take $k\le n/3$; the other thresholds are covered by the one- and two-cycle band, apart from the single transition $n=3k-1$ which has $k>n/3$. The same residual estimates in fact hold for $k\le(n+3)/3$, so the displayed proof covers this transition too.

\subsection{Expanding by occupied pole subsets}

Use $z_p=(\cos\theta+ipS)^{-1}$, $\rho_p=|z_p|$, and $F_j,D,d_j$ as in the \href{https://github.com/ProgramComputer/open-math-research/blob/main/topics/long-cycle-excedances/central-interval.md}{central interval proof}. Put $g_j=d_j$ for $j\ge k$ and zero elsewhere. Let $b_j=g_j-g_{j-1}$, so $\sum|b_j|<2/k$ and $b_j=0$ below $k$. The exact generating series is
\[
\sum_m F_m u^m=\prod_{p\,\mathrm{odd}}D(z_pu),\qquad
D(z)-1=\frac{\sum_{j\ge k} b_j z^j}{1-z}.
\]
The contribution supported on a single pole is $d_m B_m$, $B_m=\sum_p z_p^m$. Write $R_m=F_m-d_m B_m$ when $m\ge k$. For a finite set $I$ of $h\ge2$ distinct occupied poles the coefficient with total exponent $m$ is summed over $s_p\ge k$, $\sum s_p\le m$, with weight $\prod b_{s_p}$. By partial fractions of $\prod_{p\in I}(1-z_pu)^{-1}$ its remaining factor is
\[
\sum_{a\in I} \prod_{p\in I} z_p^{e_p}
\prod_{q\in I\setminus\{a\}}(1-z_q/z_a)^{-1},
\]
where $e_q=s_q$ for $q\ne a$, $e_a=s_a+m-\sum s_p$, and $\sum e_p=m$. The distinctness of the $z_p$ holds for $0<\theta<\pi$. In particular $e_p\ge k$ and some $b\in I$ has $e_b\ge m/h$.

We will bound the value for $m=n$ and $m=n+1-k\ge k$ with $mS^2\ge2/3$. The derivative bound only uses $m=n$. Write $w_p=\rho_p^{k/3-1}$. Since
\[
|1-z_q/z_a|=|q-a| S\rho_q,
\]
the absolute value of each term inside the sum over $a$, after suppressing the $b_{s_p}$ weights, is at most
\[
S^{1-h}\rho_b^{m/(3h)}\,w_a\prod_{q\ne a}\frac{w_q}{|q-a|}. \tag{PB}
\]
We took $k/3$ powers at every pole and $m/(3h)$ at $b$. A further $e_p/3$ powers remain at each pole and can be used for differentiating. We can upper bound (PB) by summing over all choices of $b$ before summing the weights. Their absolute sum for a fixed $I$ is at most $(2/k)^h$.

For the normalized derivative at $m=n$, differentiating the powers costs at most $9h/4$ in this bound. For $|p|=1$, the cost $e_p/n$ is at most one. For $|p|\ge3$, $|z_p'/z_p|\le|p|\rho_p$ and
\[
\frac{e_p}{n}|p|\rho_p^{1+e_p/3}
\le \frac{6|p|}{n(p^2-1)S^2}\le\frac94
\]
by Bernoulli's inequality with exponent $e_p/6\ge1$. Each differentiated inverse denominator has relative magnitude $\rho_q/S$, since it equals $(\cot\theta+iq)/(i(q-a))$. Thus differentiation of the denominators costs at most $(h-1)/(nS)$, giving total normalized derivative cost below $5h/2$. These bounds apply to the finite-degree coefficient for each finite set and pass to the full pole sum by local absolute convergence.

\subsection{Uniform summation of the partial fractions}

For $\lambda\ge 8$ use $(1+x)^{\lambda/2}\ge(1+(\lambda/4)x)^2$. For $|p|\ge3$, $p^2-1\ge8p^2/9$. Monotone integral comparison for the positive odd integers, then reflection, gives
\[
\sum_p\rho_p^\lambda\le2+\frac{3\pi}{4\sqrt2}\frac1{S\sqrt\lambda}.
\]
Therefore $\sum_p w_p\le2+4D_0$, and $V_h=\sum_p\rho_p^{m/(3h)}\le 2+4\sqrt h\le6h$ for $mS^2\ge2/3$ and $h\le m/k$. For each odd $a$ split the sum of $w_q/|q-a|$ at distance $2\lceil2D_0\rceil$. The near terms are bounded by the harmonic sum through $\lceil2D_0\rceil$, and the far ones by $(2+4D_0)/(4D_0)$. Hence $\sum_{q\ne a}w_q/|q-a|\le U$.

In summing (PB), the part $a=b$ costs at most $V_h U^{h-1}/(h-1)!$ after omitting $S^{1-h}$. For $b\ne a$, first sum the $h-2$ remaining poles for fixed $a,b$. The sum of $w_a/|a-b|$ over $a\ne b$ is at most $U$ for each $b$, and $\sum_b\rho_b^{m/(3h)}w_b\le V_h$. This gives at most $V_h U^{h-1}/(h-2)!$. Keeping the largest exponent at some pole is important here, since otherwise $V_h$ would depend on $n/k$. The two bounds together are $h V_h U^{h-1}/(h-1)!$.

Put $\eta=2U/(kS)\le1/1000$. Since $h^2\le 4^{h-1}$ and $h^3\le8^{h-1}$,
\[
|R_m|\le \frac{12}{k}\frac{4\eta}{1-4\eta}<\frac1{20k},\qquad
\frac{|R'_n|}{n}\le \frac{30}{k}\frac{8\eta}{1-8\eta}<\frac1{4k}. \tag{PC}
\]
All remaining poles and every number of cycles are included.

\subsection{Applying the old-root insertion comparison}

For $n\ge100$, $nS^2\ge1$ the ordinary single-pole tail obeys $\sum_{p=3,5,\ldots}p\rho_p^n<1/10$, by equation (6) of the \href{https://github.com/ProgramComputer/open-math-research/blob/main/topics/long-cycle-excedances/first-lag.md}{first two-cycle boundary argument}. Thus $|B_n-2\cos(n\theta)|<1/15$ and $|B'_n/n+2\sin(n\theta)|<1/5$. With $1/(2k-1)\le d_n\le1/k$ and (PC), the total value error is below $(1/15+1/20)/k$ and the derivative error below $(1/5+1/4)/k$. This gives the integer samples and, at every old zero, $|\cos(n\theta)|<7/60$, $|\sin(n\theta)|>19/20$, and $|F'_n|>n/(2k)$.

At a zero of $F_n$ the normalized insertion contributions are $S F'_n$ and $B_k F_m$ where $m=n+1-k$. If $mS^2\ge2/3$ and $m\ge k$, the absolute sum bound above at $\lambda=m$ yields $|B_m|\le2+3\pi\sqrt3/8< 4.1$. Thus $|F_m|<5/k$. Also
\[
\frac{|B_k|}{nS}\le \frac2{nS}+\frac{3\pi}{4\sqrt{2k}\,nS^2}
\le \frac{2}{2000}+\frac{3\pi}{4\sqrt{4000}}<\frac1{10},
\]
using $n\ge k$, $kS\ge2000$ and $nS^2\ge1$. Consequently $|B_kF_m|<nS/(2k)<S|F'_n|$.

Taking $k\ge2000\sqrt n\,[3+\log(3\sqrt{n/k})]$ suffices for (PA) throughout $nS^2\ge1$. This is an explicit effective but conservative condition. It leaves the intermediate threshold regime open.

\section{An explicit strict-interlacing range with unbounded cycle count}

For fixed $3\le k\le n$, let $E_{n,k}(t)$ count permutations of $[n]$ by excedances subject to every cycle having at least $k$ letters. Write $P_n=E_{n,k}/t$. Inversion of a permutation exchanges excedances and deficiencies. A single cycle can attain either extreme, so $P_n$ is reciprocal of degree $n-2$ with constant and leading coefficients one. \href{https://www.symmetricfunctions.com/realRootedOpenProblems.htm\#conjLongCycleExcedances}{Original conjecture} $\cdot$ \href{https://github.com/ProgramComputer/open-math-research/blob/main/topics/long-cycle-excedances/two-cycle-band.md}{Cycle weights and recurrence}

\textbf{Partial theorem.} The zeros of $P_n$ and $P_{n+1}$ are simple, negative and strictly interlace whenever $n\ge10{,}000$ and
\[
k\ge2000\sqrt n\left(3+\log\left(3\sqrt{n/k}\right)\right).
\]
The bound is effective and permits unboundedly many cycles. For example, take $n=2^{4j}$ and $k=2^{3j}$ for any integer $j\ge15$. The condition reduces to $2^j\ge2000(3+\log3+(j/2)\log2)$. Already $2^{15}>2000(5+15/2)$, and doubling the left side at each step preserves the inequality. The largest allowed cycle count is then $2^j$.

\subsection{Counting every root}

If $n\le3k-2$ this is the \href{https://github.com/ProgramComputer/open-math-research/blob/main/topics/long-cycle-excedances/two-cycle-band.md}{one- and two-cycle theorem}. In the remaining range $n\ge3k-1$, the residual index $m=n+1-k$ is at least $2n/3$ and $m\ge k$. Put $H=\lfloor7\sqrt n/20\rfloor$. The \href{https://github.com/ProgramComputer/open-math-research/blob/main/topics/long-cycle-excedances/effective-front.md}{coefficient argument} locates the $H$ smallest roots in modulus, all negative and simple, and verifies both insertion signs there. The last of these roots in modulus precedes the sample of index $H+1$.

Parameterize the negative axis by $t=-\exp(-\pi\cot\theta)$ with $0<\theta<\pi$. Then $|t|$ increases with $\theta$. On the first half of the interval take samples $\theta_j=j\pi/n$, starting with $j=H+1$. At this first sample $n\sin^2\theta_j>6/5$. All angles through $\pi/2$ consequently have $n\sin^2\theta\ge1$, $m\sin^2\theta\ge2/3$ and $k\sin\theta\ge2000(3+\log(3\sqrt{n/k}))$. The \href{https://github.com/ProgramComputer/open-math-research/blob/main/topics/long-cycle-excedances/separated-poles.md}{separated-pole estimate} gives the sample sign $(-1)^j$ for $E_{n,k}$ and controls the insertion term at every old zero in the interval. Reciprocity extends the sample signs, the old-root comparison, and the exterior-root count to the other half. For odd $n$ the central sample gap straddles $\pi/2$ and contains the reciprocal root $-1$.

The consecutive samples with indices $H+1,\dots,n-H-1$ give $n-2H-2$ sign changes. The $2H$ distinct exterior roots exhaust the remaining degree. Thus every gap contains precisely one root and every root is simple. At every old root $x$ we have
\[
P_{n+1}(x)=x(1-x)P'_n(x)+L(x),\qquad
|L(x)|<|x(1-x)P'_n(x)|.
\]
The signs of $P_{n+1}$ there are successively $-1,+1,\dots$ in order of increasing modulus. The positive constant term and the positive leading coefficient, together with $\deg P_{n+1}=n-1$, give both outside roots as well as one new root in each old gap. This proves the strict interlacing.

\section{A numerical cutoff for every minimum cycle length}

\textbf{Partial theorem.} For $n\ge10^{20}$ and every $3\le k\le n$, the nonzero zeros of the two consecutive minimum-cycle-length excedance polynomials are simple, negative and strictly interlace. This makes the earlier existential cutoff explicit. Checking the entire finite complement by enumeration is computationally impractical. The sign comparison below is for the actual next row.

Use the \href{https://github.com/ProgramComputer/open-math-research/blob/main/topics/long-cycle-excedances/effective-threshold.md}{effective large-threshold theorem} if $k\ge2000\sqrt n(3+\log(3\sqrt{n/k}))$. In the remaining range the short-cycle estimate applies. More generally, it proves strict interlacing for every $n\ge10^6$ satisfying
\[
\frac n{200k}>30\log n+20. \tag{U1}
\]
To check that the two sufficient conditions cover every threshold when $n\ge10^{20}$, observe first that if $k\le\sqrt n$ equation (U1) follows directly. Otherwise $k>\sqrt n$ and $k<2000\sqrt n(3+\log 3+(\log n)/4)$. At $n=10^{20}$, $400000(3+\log3+(\log n)/4)(30\log n+20)<10^{10}$ using $\log 3<6/5$ and $\log10<7/3$. The ratio of $\sqrt n$ to these two logarithmic factors increases for $\log n\ge40$, by logarithmic differentiation.

\subsection{A common coefficient contour}

Use $t=-e^{-s}$, $s\ge0$, $w_p=s+ip\pi$ for odd $p$, $r=|w_1|$, $S=\sin\theta=\pi/r$ and $z_p=r/w_p$. Define $T_s$ by the \href{https://github.com/ProgramComputer/open-math-research/blob/main/topics/long-cycle-excedances/fixed-threshold.md}{fixed-threshold coefficient identity}
\[
\frac{E_{j,k}(t)}{j!(1-t)^{j+1}}=[w^j]\frac{T_s(w)}{1+e^{w-s}}.
\]
The logarithm $G_s(w)$ of $T_s(w)$ is the specific polynomial with $G_s(0)=0$. Its linear term is $e^{-s}w/(1+e^{-s})$. Its terms of degree $\ell\ge2$ have absolute size at most $(|w|/r)^\ell\sum_p|z_p|^\ell/\ell$, by the Eulerian pole formula. For $s>0$ it also has the absolutely convergent expansion
\[
G_s(w)=-\sum_{a\ge1}\frac{(-e^{-s})^a}{a}\sum_{\ell=1}^{k-1}\frac{(aw)^\ell}{\ell!}.
\]
On $|w|\le R_*=r(1+1/(50k))$ we claim
\[
|G_s(w)|\le B:=6+10\log k . \tag{U2}
\]
For $s\le10$, $\sum_p|z_p|^2\le2+1/(2S^2)<8$, by telescoping $\sum_{p=3,5,\dots}1/(p^2-1)=1/4$. The linear term costs less than six, and $8e^{1/50}<10$ bounds the coefficient of the harmonic sum in the rest.

For $s>10$ put $L=\min(k-1,\lfloor s^2/16\rfloor)$ and $\delta=R_*-s\le5/s+s/(40k)$. For terms of degree at most $L$, use the exponential sum. The absolute truncated exponential, including the factor $e^{-as}$, is bounded by $e^{a\delta}\Pr(\operatorname{Pois}(aR_*)\le L)$. If $a\le2L/s$, then $a\delta\le5/8+1/20<1$, so this range contributes at most $3(1+\log k)$. For $a>2L/s$ the truncated exponential series times $e^{-aR_*}$ is bounded by the lower-tail probability $\Pr(\operatorname{Pois}(aR_*)\le aR_*/2)\le e^{-aR_*/8}$. Since $\delta/s<3/50$, the resulting sum is at most $\sum_{a\ge1}e^{-as/16}/a<2$. Any remaining degrees have $\ell\ge7$ and $\ell S^2>1/2$. By the lattice integral estimate in \href{https://github.com/ProgramComputer/open-math-research/blob/main/topics/long-cycle-excedances/separated-poles.md}{the pole argument},
\[
\sum_p|z_p|^\ell\le2+\frac{3\pi}{4\sqrt2 S\sqrt\ell}<5 .
\]
Here the integral estimate is valid for $\ell\ge4$ because $(1+x)^{\ell/4}\ge1+\ell x/4$. Thus the higher degrees cost at most $6\log k$, proving (U2).

Take $nS^2\ge1$. If $s\le10\sqrt k$ use the circle $R=r(1+1/(100k))$. Only $w_{\pm1}$ are inside, since $R_*^2-r^2<5$ whereas $|w_3|^2-r^2=8\pi^2$. The distance from the circle to each pole exceeds $n^{-2}$, using $r/(100k)\ge\pi/(100n)$ for the first pair and the squared gap for the others. If $s>10\sqrt k$, choose a radius in $[r(1+1/(100k)),R_*]$ whose distance from every pole radius is at least $n^{-2}$. This interval has length at least $\pi/(100n)$ and fewer than $10\sqrt n$ pole radii can meet it within $n^{-2}$, as $r\le\pi\sqrt n$. On both circles $|1+e^{w-s}|^{-1}\le4n^2$. Indeed, when $|\Re(w-s)|\ge n^{-2}/2$ use $|1-e^{\Re(w-s)}|$. Otherwise the distance of $\Im w$ from the nearest odd multiple of $\pi$ is at least $\sqrt3\,n^{-2}/2$, and the squared modulus is $(1-e^{\Re(w-s)})^2+4e^{\Re(w-s)}\cos^2(\Im w/2)$.

For $j=n,n+1$ use the common positive normalization
\[
Y_j=\frac{r^{j+1}E_{j,k}(t)}{2|T_s(w_1)|j!(1-t)^{j+1}} .
\]
The residues give the sum over enclosed positive odd $p$ of $\Re(T_s(w_p)z_p^{j+1})/|T_s(w_1)|$. The circle remainder in either row has magnitude at most $2n^2r\exp(2B-n/(200k))<n^{-2}$, by (U1).

\subsection{Neighboring residues}

Put $\gamma=\Im G_s(w_1)$ and $\Psi=(n+1)\theta-\gamma$. The principal pair contributes $\cos\Psi$ to $Y_n$. For $s>10\sqrt k$, at odd poles one has
\[
G_s(w_p)=-\log(1+e^{-s})-H(w_p),\qquad
H(w)=\sum_{a\ge1}\frac{e^{-aw}}a\sum_{\ell=0}^{k-1}\frac{(aw)^\ell}{\ell!}.
\]
This follows by substituting $e^{-aw_p}=(-e^{-s})^a$ in the polynomial expansion. On the vertical segments from $s$ to enclosed $w_p$ one has $\log(|w|/s)\le 7/(100k)$. Differentiating the uniformly converging series for $H$ there gives
\[
|H'(w)|\le e^{7/100}\sum_{a\ge1}e^{-as}(as)^{k-1}/(k-1)!.
\]
For $k\ge14$ the sum is at most $1/s+1/\sqrt{2\pi(k-1)}$, by summing the unimodal gamma density on a mesh of spacing $s$. Its integral is one and Stirling's lower bound gives the stated bound at its mode $k-1$. Thus $\pi |H'|\le\pi e^{7/100}(1/(10\sqrt{14})+1/\sqrt{26\pi})<12/25$. For $3\le k\le13$, $s>10\sqrt k\ge2k$ and $s>17$. Each summand is bounded by the Poisson probability of falling below half its mean, at most $e^{-as/8}$. Therefore $\pi |H'|\le\pi e^{7/100}/(e^{17/8}-1)<12/25$. Since $H(s)$ is real,
\[
|\gamma|<12/25,\qquad
|G_s(w_p)-G_s(w_1)|\le(12/25)(p+1) \quad(p\ge3).
\]

An enclosed pole has $(p^2-1)S^2\le1/(20k)$, hence $|z_p|^n\le\exp(-(49/100)(p^2-1))$. The total nonprincipal weighted error on this circle is bounded by
\[
\sum_{p=3,5,\ldots}p\,|z_p|^n\,|T_s(w_p)/T_s(w_1)|
<3e^{-2}+\sum_{p=5,7,\ldots}p e^{-(37/100)(p^2-1)}<21/50 ,
\]
where the sum on the left runs only over enclosed poles. For $p\ge5$, $(p+1)/(p^2-1)\le1/4$. The remaining sum is bounded by its first term and half the integral of the decreasing summand from $5$ to infinity. For $s\le10\sqrt k$ the nonprincipal sum is empty. Subtracting $\cos\theta\,Y_n$ from the following row multiplies each pole contribution by $z_p-\cos\theta$, whose magnitude is at most $pS$. Thus
\[
|Y_n-\cos\Psi|<1/2,\qquad |Y_{n+1}-\cos\theta\,Y_n+S\sin\Psi|< S/2. \tag{U3}
\]
These hold throughout $nS^2\ge1$. The correction $\gamma$ need only be bounded on the outer endpoint for the root count.

\subsection{A phase-crossing argument}

Let $H_0=\lfloor7\sqrt n/20\rfloor$ and $\theta_*=(H_0+5/6)\pi/n$. Then $n\sin^2\theta_*>1$, $\theta_*<1/100$ and $s_*=\pi\cot\theta_*>2\sqrt n>10\sqrt k$. These follow from $H_0+1\le(9/25)\sqrt n$, the sine and cosine Taylor bounds, and (U1). Thus $\Psi(\theta_*)<(H_0+1)\pi$ by $12/25+1/100<\pi/6$. All $H_0$ exterior roots from the \href{https://github.com/ProgramComputer/open-math-research/blob/main/topics/long-cycle-excedances/effective-front.md}{effective front proof} lie below $|t(\theta_*)|$. Indeed their moduli are bounded by $(9/8)(51/5)(H_0/(H_0+1))^n$ and
\[
n\log(1+1/H_0)-\frac n{H_0+5/6}
\ge \frac n{(H_0+1)^2}\left(\frac12-\frac{H_0+1}{6(H_0+5/6)}\right)
>\frac{33/100}{(9/25)^2}>\frac52>\log((9/8)(51/5)).
\]
Use $\pi\cot\theta_*<n/(H_0+5/6)$. At the central endpoint $s=0$, $\gamma=\pi/2$: the linear term gives this value, even powers in $G_0(i\pi)$ are real, and odd powers of degree at least three have zero coefficient by the Eulerian pole formula. Hence $\Psi(\pi/2)=n\pi/2$ exactly.

Take the first hits, after $\theta_*$, of the levels $j\pi$ for $H_0+1\le j\le\lfloor n/2\rfloor$. They occur in the indicated order by continuity. Between two consecutive hits take the last hit of the lower level before the upper one. On the intervening open interval the phase is strictly between the two levels. By (U3) the old polynomial has opposite endpoint signs. Select one zero there. Its phase gives the sign of the next polynomial by (U3), since $|\cos\Psi|<1/2$ implies $|\sin\Psi|>\sqrt3/2$. For even $n$, an early first hit of the central level is harmless: the reciprocal intervals are disjoint and the degree count below exhausts all roots. For odd $n$, the odd reciprocal polynomial $P_n$ has the additional zero $-1$ with the required next-row sign directly from (U3) at $\pi/2$. Reflecting the selected half-intervals and including the $2H_0$ exterior roots yields $n-2$ distinct negative roots for $P_n$ in both parities. Its degree is $n-2$, so each is simple. The old-root signs from (U3) and the exterior-root estimates alternate for $P_{n+1}$ after division by $t$. Reciprocity gives $P_{n+1}(1/x)=x^{-(n-1)}P_{n+1}(x)$ at the reflected roots. Its constant and leading coefficients are one and its degree is $n-1$. The intermediate value theorem gives strict interlacing. This proof needs no derivative bound or monotonicity for the phase correction.
